% GENERATED FILE. Edit canonical lessons, then run npm run generate:books.
% canonical-source-hash: fnv1a64:fdc5f1b37fb29314
% canonical-lessons: ZH-C146-baitian, ZH-C146-yeli, ZH-C146-banye, ZH-C146-xing-wake, ZH-C146-pingshi, ZH-R146-review-141-146, ZH-R146-review-final-146

\chapter{Day and Night}
\label{ch:zh-day-and-night}
\hlchaptermodality{146}

\begin{chapteropening}
\textbf{By the end of this chapter:} \emph{I can talk about the daytime, the night, the middle of the night, waking up and what I usually do.}

Complete the last lesson of chapter 146: I can talk about the daytime, the night, the middle of the night, waking up and what I usually do.
\end{chapteropening}

\section[\texorpdfstring{\zh{白天} (báitiān)}{白天 (báitiān)}]{\zh{白天} (báitiān) — daytime}

\label{lesson:ZH-C146-baitian}



\begin{quote}
Before the new one: say the Mandarin for to lie in, then the Mandarin for hurry.
\end{quote}

\subsection*{You'll want to know: \zh{白天}}
\textbf{\zh{白天}} — \emph{báitiān} — \textquotedblleft{}daytime\textquotedblright{}.

\begin{quote}
\textbf{\zh{白天这儿很吵}} \emph{báitiān zhèr hěn chǎo} — It's noisy here in the daytime.
\end{quote}

\begin{cousinweb}
\textbf{\zh{白天}} (\emph{báitiān}) is built from \zh{白} 'white' and \zh{天} 'day, sky', both inherited from Middle Chinese, \emph{baek} and \emph{then} (Guangyun, in Baxter's transcription): the white, light part of the day.
\end{cousinweb}

\begin{grammarlens}[title={What you've built}]
One more word for times of day.
\end{grammarlens}

\subsection*{Your turn}
\begin{itemize}
\raggedright
  \item \emph{Say it:} \emph{báitiān}
  \item \emph{Say it:} \emph{báitiān}, once more
  \item \emph{Say it:} \emph{báitiān zhèr hěn chǎo} — It's noisy here in the daytime.
  \item \emph{Recall:} say the Mandarin for to lie in, then the Mandarin for hurry, then say \emph{báitiān} again
\end{itemize}

\begin{tcolorbox}[breakable,colback=teal!4,colframe=teal!35!black,title={Before you move on}]
What does \zh{白天} mean? (\textquotedblleft{}Daytime\textquotedblright{}.) Say it once more.
\end{tcolorbox}
\section[yèli]{yèli — at night}

\label{lesson:ZH-C146-yeli}



\begin{quote}
Before the new one: say the Mandarin for hurry, then the Mandarin for daytime.
\end{quote}

\subsection*{You'll want to know: yèli}
\emph{yèli} — \textquotedblleft{}at night\textquotedblright{}.

\begin{quote}
\emph{yèli hěn lěng} — It's cold at night.
\end{quote}

Its characters come later; for now it is a word for the ear.

\begin{cousinweb}
\emph{yèli} is built from \emph{yè} 'night' and \zh{里} 'inside', both inherited from Middle Chinese, \emph{yaeH} and \emph{liX} (Guangyun, in Baxter's transcription); the dictionary glosses the first 'dusk' and 'a time to rest'.
\end{cousinweb}

\begin{grammarlens}[title={What you've built}]
One more word for times of day.
\end{grammarlens}

\subsection*{Your turn}
\begin{itemize}
\raggedright
  \item \emph{Say it:} \emph{yèli}
  \item \emph{Say it:} \emph{yèli}, once more
  \item \emph{Say it:} \emph{yèli hěn lěng} — It's cold at night.
  \item \emph{Recall:} say the Mandarin for hurry, then the Mandarin for daytime, then say \emph{yèli} again
\end{itemize}

\begin{tcolorbox}[breakable,colback=teal!4,colframe=teal!35!black,title={Before you move on}]
What does yèli mean? (\textquotedblleft{}At night\textquotedblright{}.) Say it once more.
\end{tcolorbox}
\section[bànyè]{bànyè — the middle of the night}

\label{lesson:ZH-C146-banye}



\begin{quote}
Before the new one: say the Mandarin for daytime, then the Mandarin for at night.
\end{quote}

\subsection*{You'll want to know: bànyè}
\emph{bànyè} — \textquotedblleft{}the middle of the night\textquotedblright{}.

\begin{quote}
\emph{wǒ bànyè è le} — I got hungry in the middle of the night.
\end{quote}

Its characters come later; for now it is a word for the ear.

\begin{cousinweb}
\emph{bànyè} is built from \zh{半} 'half' and \emph{yè} 'night', both inherited from Middle Chinese, \emph{panH} and \emph{yaeH} (Guangyun, in Baxter's transcription).
\end{cousinweb}

\begin{grammarlens}[title={What you've built}]
One more word for times of day.
\end{grammarlens}

\subsection*{Your turn}
\begin{itemize}
\raggedright
  \item \emph{Say it:} \emph{bànyè}
  \item \emph{Say it:} \emph{bànyè}, once more
  \item \emph{Say it:} \emph{wǒ bànyè è le} — I got hungry in the middle of the night.
  \item \emph{Recall:} say the Mandarin for daytime, then the Mandarin for at night, then say \emph{bànyè} again
\end{itemize}

\begin{tcolorbox}[breakable,colback=teal!4,colframe=teal!35!black,title={Before you move on}]
What does bànyè mean? (\textquotedblleft{}The middle of the night\textquotedblright{}.) Say it once more.
\end{tcolorbox}
\section[xǐng]{xǐng — to wake up}

\label{lesson:ZH-C146-xing-wake}



\begin{quote}
Before the new one: say the Mandarin for at night, then the Mandarin for the middle of the night.
\end{quote}

\subsection*{You'll want to know: xǐng}
\emph{xǐng} — \textquotedblleft{}to wake up\textquotedblright{}.

\begin{quote}
\emph{wǒ bànyè xǐng le yīnwèi wàibian hěn chǎo} — I woke up in the middle of the night because it was noisy outside.
\end{quote}

Its characters come later; for now it is a word for the ear.

\begin{cousinweb}
\emph{xǐng} is inherited from Middle Chinese \emph{sengX} (Guangyun, in Baxter's transcription), glossed 'drunkenness passing off', that is, sobering up. Waking from sleep is not given there, so how the word came to it is not settled.
\end{cousinweb}

\begin{grammarlens}[title={What you've built}]
One more word for times of day.
\end{grammarlens}

\subsection*{Your turn}
\begin{itemize}
\raggedright
  \item \emph{Say it:} \emph{xǐng}
  \item \emph{Say it:} \emph{xǐng}, once more
  \item \emph{Say it:} \emph{wǒ bànyè xǐng le yīnwèi wàibian hěn chǎo} — I woke up in the middle of the night because it was noisy outside.
  \item \emph{Recall:} say the Mandarin for at night, then the Mandarin for the middle of the night, then say \emph{xǐng} again
\end{itemize}

\begin{tcolorbox}[breakable,colback=teal!4,colframe=teal!35!black,title={Before you move on}]
What does xǐng mean? (\textquotedblleft{}To wake up\textquotedblright{}.) Say it once more.
\end{tcolorbox}
\section[píngshí]{píngshí — usually; ordinarily}

\label{lesson:ZH-C146-pingshi}



\begin{quote}
Before the new one: say the Mandarin for the middle of the night, then the Mandarin for to wake up.
\end{quote}

\subsection*{You'll want to know: píngshí}
\emph{píngshí} — \textquotedblleft{}usually; ordinarily\textquotedblright{}.

\begin{quote}
\emph{wǒ píngshí qī diǎn qǐchuáng} — I usually get up at seven.
\end{quote}

Its characters come later; for now it is a word for the ear.

\begin{cousinweb}
\emph{píngshí} is built from \emph{píng} 'level, even' and \zh{时} 'time', both inherited from Middle Chinese, \emph{bjaeng} and \emph{dzyi} (Guangyun, in Baxter's transcription): the level, ordinary times.
\end{cousinweb}

\begin{grammarlens}[title={What you've built}]
One more word for times of day.
\end{grammarlens}

\subsection*{Your turn}
\begin{itemize}
\raggedright
  \item \emph{Say it:} \emph{píngshí}
  \item \emph{Say it:} \emph{píngshí}, once more
  \item \emph{Say it:} \emph{wǒ píngshí qī diǎn qǐchuáng} — I usually get up at seven.
  \item \emph{Recall:} say the Mandarin for the middle of the night, then the Mandarin for to wake up, then say \emph{píngshí} again
\end{itemize}

\begin{tcolorbox}[breakable,colback=teal!4,colframe=teal!35!black,title={Before you move on}]
What does píngshí mean? (\textquotedblleft{}Usually; ordinarily\textquotedblright{}.) Say it once more.
\end{tcolorbox}
\section[(dialogue)]{Looking back: from Travel and Fun to Day and Night}

\label{lesson:ZH-R146-review-141-146}



\begin{quote}
Say the Mandarin for to travel and a game.
\end{quote}

\subsection*{Your turn}
Say these aloud:

\begin{itemize}
\raggedright
  \item to travel, a game, to stroll the shops, to take a walk, to go fishing
  \item to mop the floor, rubbish, to dry in the sun, to sort out, to water the plants
  \item second (of time), minute (a length of time), to be short of, punctual, a long while
  \item important, easy, simple, special, to find (that)
  \item a reason, to run behind schedule, alarm clock, to lie in, hurry
  \item daytime, at night, the middle of the night, to wake up, ordinarily
\end{itemize}

\begin{tcolorbox}[breakable,colback=teal!4,colframe=teal!35!black,title={Before you move on}]
To travel? (\emph{lǚyóu}.) A game? (\emph{yóuxì}.) To stroll the shops? (\emph{guàngjiē}.) To take a walk? (\emph{sànbù}.) To go fishing? (\emph{diàoyú}.) To mop the floor? (\emph{tuōdì}.) Rubbish? (\emph{lājī}.) To dry in the sun? (\emph{shài}.) To sort out? (\emph{zhěnglǐ}.) To water the plants? (\emph{jiāohuā}.) Second (of time)? (\emph{miǎo}.) Minute (a length of time)? (\emph{fēnzhōng}.) To be short of? (\emph{chà}.) Punctual? (\emph{zhǔnshí}.) A long while? (\textbf{\zh{半天}}, \emph{bàntiān}.) Important? (\textbf{\zh{重要}}, \emph{zhòngyào}.) Easy? (\emph{róngyì}.) Simple? (\emph{jiǎndān}.) Special? (\emph{tèbié}.) To find (that)? (\emph{juéde}.) A reason? (\emph{yuányīn}.) To run behind schedule? (\textbf{\zh{晚点}}, \emph{wǎndiǎn}.) Alarm clock? (\emph{nàozhōng}.) To lie in? (\emph{shuì lǎnjiào}.) Hurry? (\emph{gǎnkuài}.) Daytime? (\textbf{\zh{白天}}, \emph{báitiān}.) At night? (\emph{yèli}.) The middle of the night? (\emph{bànyè}.) To wake up? (\emph{xǐng}.) Ordinarily? (\emph{píngshí}.)
\end{tcolorbox}
\section[(dialogue)]{Day and Night — once more}

\label{lesson:ZH-R146-review-final-146}



\begin{quote}
Say the Mandarin for daytime and at night.
\end{quote}

\subsection*{Your turn}
\begin{itemize}
\raggedright
  \item \emph{Say it:} daytime, at night, the middle of the night, to wake up, ordinarily
\end{itemize}

\begin{tcolorbox}[breakable,colback=teal!4,colframe=teal!35!black,title={Before you move on}]
Daytime? (\textbf{\zh{白天}}, \emph{báitiān}.) At night? (\emph{yèli}.) The middle of the night? (\emph{bànyè}.) To wake up? (\emph{xǐng}.) Ordinarily? (\emph{píngshí}.)
\end{tcolorbox}
